\documentclass[11pt,b5paper]{scrbook}
\usepackage{scrlayer-scrpage}

\KOMAoption{chapterprefix}{true}
\KOMAoption{parskip}{half}

\usepackage[svgnames]{xcolor}
\usepackage[english]{babel}
\usepackage{textcomp}
\usepackage{chapterthumb}

\AddLayersToPageStyle{@everystyle@}{chapterthumb}
\addtokomafont{chapterthumb}{\bfseries}
\usepackage{setspace}
\onehalfspacing
\usepackage{tikz}
\usepackage{graphicx}
\usetikzlibrary{graphs, positioning, calc, fit, arrows.meta}
\usepackage{tocbasic}
\usepackage[avantgarde]{quotchap}
\usepackage{lipsum}
\usepackage{geometry}
\usepackage[most]{tcolorbox}
\usepackage[chapter]{minted}
\usepackage{caption}
%\usemintedstyle{borland}

% colored DiegOS logo
\newcommand{\myOS}{{\color{SteelBlue}\textbf{DiegOS}}}

%use numbers for nested enums
\renewcommand{\labelenumii}{\arabic{enumi}.\arabic{enumii}}
\renewcommand{\labelenumiii}{\arabic{enumi}.\arabic{enumii}.\arabic{enumiii}}
\renewcommand{\labelenumiv}{\arabic{enumi}.\arabic{enumii}.\arabic{enumiii}.\arabic{enumiv}}

% Custom command to draw a stop sign
\newcommand{\stopsign}[1][1.0]{%
  \begin{tikzpicture}[baseline=-0.5ex, scale=#1]
    % Outer red octagon
    \fill[red] (22.5:1) -- (67.5:1) -- (112.5:1) -- (157.5:1)
            -- (202.5:1) -- (247.5:1) -- (292.5:1) -- (337.5:1) -- cycle;
    % Inner white border line
    \draw[white, line width=1.5pt, double=red, double distance=0.5pt]
         (22.5:0.92) -- (67.5:0.92) -- (112.5:0.92) -- (157.5:0.92)
      -- (202.5:0.92) -- (247.5:0.92) -- (292.5:0.92) -- (337.5:0.92) -- cycle;
    % Text in the center
    \node[white, font=\sffamily\bfseries] at (0,0) {STOP};
  \end{tikzpicture}%
}

\newtcolorbox{warning}[1][]{
enhanced,before skip balanced=2mm,after skip balanced=3mm,
boxrule=0.4pt,left=8mm,right=2mm,top=1mm,bottom=1mm,
colback=yellow!50,
colframe=yellow!20!black,
sharp corners,rounded corners=southeast,arc is angular,arc=3mm,
underlay={%
	\path[fill=tcbcolback!80!black] ([yshift=3mm]interior.south east)--++(-0.4,-0.1)--++(0.1,-0.2);
	\path[draw=tcbcolframe,shorten <=-0.05mm,shorten >=-0.05mm] ([yshift=3mm]interior.south east)--++(-0.4,-0.1)--++(0.1,-0.2);
	\path[fill=yellow!50!black,draw=none] (interior.south west) rectangle node[white]{\Huge\bfseries !} ([xshift=6mm]interior.north west);
	},drop fuzzy shadow,#1
}

\newtcolorbox{danger}[1][]{
enhanced,before skip balanced=2mm,after skip balanced=3mm,
boxrule=0.4pt,left=17mm,right=2mm,top=1mm,bottom=1mm,
colback=red!50,
colframe=red!20!black,
sharp corners,rounded corners=southeast,arc is angular,arc=3mm,
underlay={%
	\path[fill=tcbcolback!80!black] ([yshift=3mm]interior.south east)--++(-0.4,-0.1)--++(0.1,-0.2);
	\path[draw=tcbcolframe,shorten <=-0.05mm,shorten >=-0.05mm] ([yshift=3mm]interior.south east)--++(-0.4,-0.1)--++(0.1,-0.2);
	\path[fill=red!50!black,draw=none] (interior.south west) rectangle node[white]{\stopsign[0.75]} ([xshift=15mm]interior.north west);
	},drop fuzzy shadow,#1
}

%\lofoot{DiegOS Operating System}
%\refoot{DiegOS Operating System}

\begin{document}
\frontmatter
\subject{Documentation}
\title{{\color{SteelBlue}DiegOS}}
\subtitle{Design and Implementation of a minimalistic Operating System}
\author{Diego Gallizioli}
\date{\today}
\uppertitleback{
DiegOS\\
Design and Implementation of a minimalistic Operating System\\
by Diego Gallizioli

Copyright \textcopyright \today \\{\KOMAScriptVersion}\\
}
\lowertitleback{This work was created with the {\LaTeX} typesetting system.\\
The layout was set with {\KOMAScript}.\\
The source text was written almost entirely using {\textit{gvim}}.\\

All applications and tools were developed on an Intel workstation running Slackware Linux, ultimately the workstation is running Fedora Linux.

While every precaution has been taken in the preparation of this book, the author
assume no responsibility for errors or omissions, or for damages resulting from the use of the
information contained herein.}

\dedication{This work is dedicated to my family Monica, Davide, Matteo and Anna, for letting me sit in front of my computer late at night.}

\maketitle
\tableofcontents
\listoflistings
\mainmatter
\chapter{Introduction}
\dictum[Darth Vader]{
I find your lack of faith disturbing.
}
\newpage
\section*{Inheritance and other inspirations}
I approached the world of computing when I was 8 years old.

Back at that time information technology was seen as the future and was
represented as a kind of magic, something that could not be easily understood 
or taught in schools but rather was to be discovered.
Lots of movies used computer science as a subject to create stunning stories
where young guys (just like me!) experienced unbelievable adventures after
discovering new hidden powers while hacking stuff, not to mention the fact
that most of the offices seen in movies were filled with computers, no matter
what they could be used for, they were giving answers: finding clues, finding
people, finding informations.
Remeber \textit{tron} or \textit{war games} ? Here it goes.

My father bought a Commodore C64 after an intensive and extensive search throughout
most of the electronics shops
available, probably inspired by the passion I demonstrated for computers.

He had a target --- having me and my brother understand what a computer could
do --- but it was only partially met: I went on to become an engineer,
my brother chose a different path.\footnote{He is a doctor now.}

I still remember my joy and my interest when I found out there was more than
a BASIC interpreter in that machine.
I wondered how developers could code games or applications by using such a
simple language, so I tried myself writing code for simple programs writing
text to screen; I was constantly running out of memory, or waiting eons for the data cassette
to read in data, and getting frustrated as I could not understand "the magic"
behind the scene.
But then one day a friend of mine paid a visit to me and lent me a fabulous
gadget: a cartridge running a Pascal compiler, along with a technical manual
covering the C64 internals and the assembler for the processor and the sprites
technology... and lots more.

Then I turned 15 when I asked for a more powerful computer: my intense desire was
to explore the unknown universe of the then innovative AMIGA
platform, but to have access to that machine, I'd had to explain to my father my
reasons and also provide him with a good reason to spend his money. I failed,
because to his eyes the Commodore machines were merely gaming machines, not serious
pieces of hardware, so I had to switch to a backup plan.

I succeded in replacing my beloved C64 with an Intel 8086-based personal computer
running MS-DOS version 3.3 (\textbf{yes}, version 3.3 !)
and from then on I got more and more interested in two aspects of computing:
transmitting data and having software interact with hardware.
I approached Intel assembler, way more complicated than 6510, and then structured
procedural programming by learning Pascal.
Next came the C language, then Delphi, then 3D programming, in a never ending
spiral of learning, trying, and looking for new challenges.

While attending university I fell in love with operating systems.
This interest eventually led me to try to write an operating system by myself.
Developing a full-blown operating system as a spare time activity proved to be
really challenging. There are way too much
OS in the world let's be honest, it is not worth the effort to add a new one
just for the sake of it.
As a student I simply put aside this project, and focused on more interesting
distractions.\footnote{You might guess what could possibly distract a 25 years 
old man from spending nights in front of a computer.}
But, as a professional software developer I approached embedded systems and I found 
an interesting echosystem to work on.
This is why I resumed my late OS project and wrote {\myOS}.

From my personal standpoint the main goal of an Operating System is to offer a 
friendly and comfortable environment for both programmers, users, and administrators,
or in other words, offer an effective interface to the underlying hardware.
Scoring this goal is not easy: since the very first Operating System inception, a huge number of engineers 
and scientists worked to develope concepts, ideas, architectures, interfaces, and documenting 
their work to produce what we know today as Microsoft Windows, Linux, Unix, or Apple Inc. macOS.
All these people proved to be smarter than the Author of this book, no need to
argue upon it.

Up today I had the chance and the privilege to work on a number of Operating Systems,
some of them with peculiar features, some of them with legacy constraints, most of
them with a clear design.
{\myOS} was born after years of careful and intensive meditation; it has no peculiarities,
a KISS \footnote{The acronim stands for Keep It Simple, Stupid.} approach, no legacy constraints, and above all a humble
target: do the right thing, the easy way, the fastest way, avoiding common errors and nuisances.

Development of {\myOS} took a lot of my spare time, still I did not conceive all the contents of this book by myself;
the most inspiring source of wisdom at the heart of
{\myOS} is \textbf{Minix}.
Minix is definitively a fitness room, furnished with all the whistles and bells and tools you could ever dream of.
This \textbf{\emph{does not imply}} that Minix is the best Operating System out there, 
it simply means I consider it a delightful piece of code.

Besides, the most inspiring source to improve my programming skills was
Michael Abrash's \emph{Graphics Programming Black Book};
all I have to say about it is, zenning is a very long journey, and you'd better take your time and read it.

Last but not least, writing software requires patience. Resilience to failure. Acceptance of personal limits.
Acknowledge of success when new code runs the way it should.
Fighting anger and frustration when things just don't go.
Tolerance with the "Muggles" not willing to see the magic running under the hood.
Finally, writing software requires understanding of the true nature of software: vacuity.
Because all software of any kind under any form is
essentially an abstract algorithm performing elementary actions faster than
our mind could.
The creator of all software \textbf{is} our mind, but \textbf{no software can replace} our mind.\footnote{Now in 2025 it seems that \textbf{AI} will prove me wrong.}

I leave to the reader an abstract taken from \textit{The Heart Sutra}, my favourite lecture when in need of wiseness.

I hope you will enjoy my software.

\emph{" Form is empty, emptiness is form, Emptiness is not other than form, form is also not other than emptiness.
Likewise, sensation, discrimination, conditioning, and awareness are empty.}

\emph{In this way, Sariputra, all things are emptiness; they are without defining characteristics;
they are not born, they do not cease, they are not defiled, they are not undefiled.
They have no increase, they have no decrease.}

\emph{Therefore, Sariputra, in emptiness there is no form, no sensation, no discrimination, no conditioning, and no awareness.
There is no eye, no ear, no nose, no tongue, no body, no mind.
There is no form, no sound, no smell, no taste, no texture, no phenomenon.
There is no eye-element and so on up to no mind-element and also up to no element of mental awareness.
There is no ignorance and no elimination of ignorance and so on up to no aging and death and no elimination of aging and death.
Likewise, there is no suffering, origin, cessation, or path; there is no wisdom, no attainment, and even no non-attainment.''}
\vfill

\chapter{The Architecture}
\dictum[C.A.R. Hoare, The 1980 ACM Turing Award Lecture]{
There are two ways of constructing a software design:
One way is to make it so simple that there are obviously
no deficiencies and the other way is to make it so complicated
that there are no obvious deficiencies.
}
\newpage
\section{Design}
{\myOS} has been designed to be a simple application running on top of bare metal.
It is compiled and linked as a single binary image, a single process application able to run 
concurrent threads of execution by means of a cooperative scheduler.
The software is partitioned in layers with defined dependencies between modules;
{\myOS} does not support multiple processors,  multiple cores, or multiple hardware 
threads of execution; it does not offer memory protection or swappable memory,
and does not offer pre-emption.
This kind of design is not innovative, it resembles most embedded OS available
online, or what are commonly known as realtime OS.

The hardware abstration is implemented with a well defined virtual interface
architecture commonly known as \textbf{HAL} or Hardware Abstraction Layer.

{\myOS} is designed to be configurable, and the
first configuration step is
the definition of a set of options to be used at compile time to instruct the code and
the compiler on what to do with very specific aspects of the system, like
context switching, interrupt handling, size of memory, size of buffers,
maximum number of threads, instruction set to be used.

\subsection{Overview}
Using a bottom-up perspective the system is split into 3 layers with foundations consisting of
the \textbf{HAL} for the kernel and the hardware for \textbf{Drivers}.
Drivers are logically next to the kernel and the HAL because they access hardware directly
(this is shared with the HAL) and they are a part of the kernel (same privilege but no abstraction).
\hfill

\begin{tcolorbox}[enhanced,adjusted title=Hardware,drop fuzzy shadow]
    \begin{tcbitemize}[
        raster columns=3,
	raster rows=4,
        raster column skip=5mm,
        raster row skip=5mm,
	raster equal height=rows,
        halign=center]
        
        \tcbitem[
		raster multicolumn=3,
		colback=lightgray!90!green,
		colframe=green!50!black,
		enhanced,
		adjusted title=Layer 3,
		halign title=center,
		drop fuzzy shadow]
            Applications
        \tcbitem[
		raster multicolumn=3,
		colback=lightgray!90!yellow,
		colframe=yellow!80!darkgray,
		enhanced,
		adjusted title=Layer 2,
		halign title=center,
		drop fuzzy shadow]
            Libraries
	\tcbitem[raster multirow=2,raster multicolumn=2,blankest]
		\begin{tcbitemize}[raster columns=2,raster rows=2]
	        \tcbitem[
			raster multicolumn=2,
			colback=lightgray!90!cyan,
			colframe=cyan!40!darkgray,
			enhanced,
			adjusted title=Layer 1,
			halign title=center,
			drop fuzzy shadow]
	            Kernel
		 \tcbitem[
			raster multicolumn=2,
			colback=lightgray!90!red,
			colframe=red!60!gray,
			enhanced,
			adjusted title=HAL,
			halign title=center,
			drop fuzzy shadow]
	            Hardware Abstraction Layer
		\end{tcbitemize}

        	\tcbitem[
			colback=lightgray!90!violet,
			colframe=violet!50!black,
			enhanced,
			adjusted title=Layer 1a,
			halign title=center,
			drop fuzzy shadow]
	            Drivers
   \end{tcbitemize}
\end{tcolorbox}
\hfill

At the foundations of {\myOS} there is the {\color{Red}\textbf{HAL}} which is a virtualization of
\marginline{\footnotesize{the system foundations}}
common functions or operations to be performed by hardware or by platform-specific instructions.
This part of the software is responsible for the correct execution of a context switch,
access to memory-mapped registers, access to the Memory Management Unit or the Memory Protection Unit,
BIOS management and parameters access, possibly BIOS routines execution
\footnote{This specific capability may or may not be available, it depends on the processor and
its execution mode when entering BIOS code.},
handling of boot parameters, and a lot more.

Side-by-side with the HAL there are {\color{Violet}\textbf{Drivers}} with a well-defined interface used to access hardware
devices. Drivers are sensitive and prone to fatal errors and need to mess with
interrupts. They are logically side-by-side with the kernel as highest privileges are required to
handle hardware resources and no other part of the system is allowed to run drivers --- only the kernel
can create an instance of a driver and start it.

The starting layer is the {\color{Cyan}\textbf{Kernel}},
\marginline{\footnotesize{1\textsuperscript{st} layer: the Kernel}}
composed of several items like a mutex manager to coordinate
concurrency between threads, an alarm manager to implement timers, an event manager
to exchange messages between threads, barriers, spin locks, \textbf{I/O} support for drivers.
The inner part of the core is all about thread management and is not directly exposed
to other modules, a specific kernel API can be used to access those components
from applications.

The middle layer is occupied by {\color{Yellow}\textbf{Libraries}}, any kind of API or utility
\marginline{\footnotesize{2\textsuperscript{nd} layer: the Libraries}}
which can be used by both the kernel and the applications.
They are logically above the first layer as they implement an abstraction for algorithms
 and data structures to use raw memory, files or network connections.

In {\myOS} both the kernel and the applications can
use public functions exposed by the library layer.
In a more restrictive design the kernel would never share code or data with
applications, and this was the starting point also for {\myOS}, until it became
clear that code segregation was bloating the kernel and memory footprint was growing due to duplication
in kernel and library codebases.

In order to keep the binary size as small as possible a compromise is put in
place, which is, to share code --- and not data --- between kernel and
applications.

The top layer is occupied by {\color{Green}\textbf{Applications}} which are the actual
\marginline{\footnotesize{3\textsuperscript{rd} layer: the Applications}}
portion of the code written by developers to do something useful. Although
\textbf{applications} recalls user space processes there is no concept of user and 
kernel space in {\myOS}, applications are just code linked to the OS but using the 
APIs offered by the kernel and the libraries.
Threads, processes and applications are equivalent in {\myOS}.
Still, applications are not allowed to call the kernel's internals directly.
The kernel is an indipendent item and need to be isolated as much as possible.

The following sections provide further details for each layer.

\subsection{The Hardware Abstraction Layer}
\label{subsec:hal}
Software need hardware to run; a lot of software components need to interact 
with some piece of hardware which is able to offer features not achievable in
other ways.
For instance, a software alarm need a timer to measure time and understand when
the alarm is expired; a timer need to count seconds with the best resolution possible,
and the reference source of this "ticking" is usually a crystal oscillator.
Measuring the crystal signal is performed in hardware as it is an analog signal which
is sampled, quantized, and accumulated in a counting register.
The output of hardware is a measure of time, and the software will
use it to track time passing by; but from a logical point of view the
function of time accounting \textbf{does not} depend on hardware as it is purely abstract,
seconds and days are the same quantities indipendently from the hardware they are computed on.
.\footnote{Hardware timers may have different resolutions and different
clock sources, nevertheless they perform the same task with the same logic.}
Software need to be informed when a peripheral (hardware) changes state, or
fails, or has data ready to be handled.
This notification is achieved by interrupting the processor --- so interrupting software
--- and putting both hardware and software in the condition of handling the
event.
In order to understand which is the source of the interruption the hardware provides
an \textit{interrupt} system,
offering informations on the origin of the interruption, the type of
interruption and its actual priority.
Hardware implementation and its interface toward software can vary on different platforms,
but from a logical point of
view the operation to be performed is the same: software need to know which
device asked for attention.

The \textbf{HAL} layer can be used to have a common interface between
\textit{software} and \textit{hardware functionality}; the abstraction created
by this layer has to be generic enough to cover most hardware features, while
the same should be invisible to the caller.
In few words, the software invoking the Hardware Abstraction Layer expects an
output from the API which is consistent and identical on different platforms;
the same software calling the same \textbf{HAL} interface callback will output the same
results independently from the underlying hardware.

A good example of what an \textbf{HAL} layer is capable of is interrupt handling;
interrupts generated by external devices can be processed by
dedicated peripherals\footnote{Devices managing multiple interrupt sources
are devised as PIC Programmable Interrupt Controllers.} and sorted in priority
before being forwarded to the processor. Many processors embed a dedicated device --- known
as \textit{PIC} --- so that they do not need additional devices and can offer a wider set of features
as the interrupt controller is actually on-board and can interact with other
parts of the processor.

From the \textbf{HAL} point of view there is no difference in having an
embedded PIC or an external PIC: the software interface will accept the
correct amount of interrupts, the correct amount of priorities, and will let
the software map the hardware interrupt lines to their logical counterpart in
software. When an interrupt is triggered, the \textbf{HAL} layer will call a
specific routine in charge of servicing the interrupt.
The operating system or any driver will then just need to map interrupts,
register callbacks, and run.
There is a clear tradeoff between performance and abstraction when making use
of a \textbf{HAL} layer: abstraction needs software, software needs execution time,
and execution time raises delay.
In reverse the ability to use multiple hardware with the same interfaces increase
the test surface of the software which increase stability in the long period.

\subsection{Layer 1a: The drivers}
\textbf{Drivers} are a very important part of the kernel, they are allowed to access
internals of the kernel \textbf{not} exposed to applications, and usually
implement \textit{interrupt routines} which requires access to sensitive
resources offered by the processor or the platform or the kernel itself.
It is not possible to implement a single common interface for all \textbf{drivers}
because the hardware they are supposed to use show different behaviors and require
different logics when operated; likewise the same hardware might be driven by
different drivers with different behaviors.\\
As an example a driver implementing a \texttt{TTY} might use a graphic card to print
text in lines on video; a driver implementing graphics might use the same graphic card to plot
lines and circles.\\
{\myOS} interfaces have been conceived to share ("inherit") a common
set of properties, adding specializations when needed.\footnote{Just like objects inheritance
it is possible to extend a function calls structure with another function calls structure.}

\subsection{Layer 1: The kernel}
Going a bit deeper into details, the \textbf{kernel} is composed of 4 specific
items: the \textbf{scheduler}, the \textbf{system clock manager}, the
\textbf{threads synchronization APIs} and the \textbf{kernel interface}.

The \textbf{scheduler}
\marginline{\footnotesize{scheduling operations}}
is in charge of selecting a thread among all
available threads in the system to run it; the algorithm implemented to make the selection
is a linear search in a priority list without pre-emption or substitution.
Only the running thread can perform a context switch because pre-emption is not allowed, this means
that the running thread is the only task that let other threads run by suspending itself
or waiting on events generated by other processes or by hardware.

A context switch is a voluntary action performed by the running thread and enabled by specific calls to the kernel.
Context switching \textbf{must not} happen in interrupt context, and no interrupt
routine is allowed to call any scheduler function; doing so will surely result in
a memory corruption.

As no other task 
\marginline{\footnotesize{the idle thread}}
than the running thread is able to stop or suspend itself there need to be at least one running thread
in the system at any time to allow new threads to be created and then put in operation.
This thread is usually known as the \tcboxverb[colframe=white]{idle thread} because it runs when the system is idle and the system is not busy with users' tasks.

The \tcboxverb[colframe=white]{idle thread} is fundamental for the Operating System as well as it is for the hardware management:
when the system goes idle there are no pending tasks to be executed, this is the perfect time to lower power consumption cutting off
the processor frequency, putting peripherals in power-save mode, or putting the processor in power-save mode.

The running thread 
\marginline{\footnotesize{thread execution}}
is the highest priority thread in ready state, not blocked nor waiting for an I/O operation to complete.
The \textbf{scheduler} algorithm follows a cooperative model: the running thread will not be substituted until the code in execution
\textbf{voluntary} suspends the thread, or put it in a wait state for I/O or waiting for a mutex, a
semaphore, a barrier or an event; once the running thread has changed state, the
\textbf{scheduler} will select another thread to become the running one or --- which is perfectly acceptable ---
select the same thread again if it is the first ready thread found in the priority list.

The scheduler follows a strict priority selection: the running thread will never
switch to a lower priority thread \footnote{This is a specific scheduler policy: the selection of the
running thread does not step priority down under normal conditions.}
unless the following two conditions are true:

\hfill
\begin{tcolorbox}[enhanced, colframe=CadetBlue!80!White, leftrule=3mm, top=5mm, bottom=5mm, halign=center, drop fuzzy shadow]
\begin{itemize}
\item There are no ready threads with a higher priority than the running thread
\item There are no ready threads with the same priority of the running thread
\end{itemize}
\end{tcolorbox}
\hfill

The \texttt{kernel API} offers two suspension points --- with a slightly different logic but still
honoring these rules --- in the form of function calls to let the running thread switch context:

\hfill
\begin{tcolorbox}[enhanced, colframe=YellowGreen!80!White, leftrule=3mm, top=5mm, bottom=5mm, halign=center, drop fuzzy shadow]
\begin{itemize}
\item \texttt{thread\_suspend()} suspends unconditionally the running thread.
\item \texttt{thread\_may\_suspend()} supends the running thread only if there are no higher priority threads ready to run.
\end{itemize}
\end{tcolorbox}
\hfill

As already mentioned when no user threads are running the scheduler will fall back to the 
\texttt{idle thread}; this is a special thread used to occupy the processor when there is no
specific task to be completed.
The \texttt{idle thread} performs accounting and data gathering to build a statistics dataset
of memory occupation, free stacks memory, running time and possibly many others.
Still, there is one operation performed by this thread --- most important above all --- which is of high interest
these days: power saving function, part of the \texttt{Hardware Abstraction Layer}
and to be implemented for each processor; power saving is one of the most
important functionality of the kernel, as the processor will spare power, save
batteries, and require less cooling.

The \textbf{system clock manager} is a sensitive piece of code dealing with
\marginline{\footnotesize{clock management}}
the system clock and the system timings in general.
Its responsibility is to track time --- wall clock --- and allow threads to wait for a specific amount
of time.\\
It will initialize the clock devices\footnote{Version 1.x of {\myOS} supports
one clock source only \texttt{/dev/clk0}} and perform accounting of system
ticks\footnote{System ticks are implemented as a monotonic fixed rate counter},
boot time\footnote{The system boot time is a configurable time point marking a date
since 1\textsuperscript{st} January 2015 GMT}, and the clock \textit{period}.

A clock device is simply a counter; some devices count upwards from zero to higher values,
some count backwards to zero. In any case the value of the counter is called the \textit{period}.
Once the \textit{period} is expired and the counter reaches its terminal value an interrupt is raised
and then the counter may restart counting or may wait to be restarted;
the two different approaches are usually referred as \textit{periodic} and \textit{single-shot}.

The counter may or may not raise and interrupt,
\marginline{\footnotesize{system ticks}}
in the latter case we have a free-running counter which can
be polled to feed a variable in memory accumulating the time passing by.
{\myOS} compute \textit{system ticks} by accumulating \textit{periods} as configured into the clock device;
the resolution of the system clock is dependent on the resolution of the hardware.
The system clock can be accessed by the application layer with different resolutions

\hfill
\begin{tcolorbox}[enhanced, colframe=LightSalmon!60!White, coltitle=Black, adjusted title=Time measurement, top=5mm, bottom=5mm, halign=center, drop fuzzy shadow]
\begin{center}
    \begin{tabular}{| c | l |}
    \hline
    System Clock & Elapsed time in milliseconds \\ \hline
    Up time & Adjusted elapsed time since boot \\ \hline
    System ticks & Elapsed time in HZ ticks per second \\
    \hline
    \end{tabular}
\end{center}
\end{tcolorbox}
\hfill

System ticks are computed dividing the elapsed time by a specific frequency, usually configured at compile time
declaring a \texttt{HZ} definition\footnote{Most common HZ values are 100, 1000 and 1024, with 100 the default
choice for {\myOS}}.
{\myOS} use a millisecond resolution for its timers while hardware can be configured to have any resolution,
the software will account for milliseconds by converting the accumulated elapsed time at run time.
One note about having monotonic clocks and not real-time clocks: it might seem an oversimplification, but experience (and history!)
teach a different lesson: real-time clocks can confuse the status of a system by moving timepoint up and down the timeline.
This poses a major issue for threads waiting for a specific timepoint, in fact, an unbounded operation on a timepoint
(i.e. mixing signed and unsigned additions) might set a timepoint in the past or very far in the future, in both cases
the wait operation will fail its original goal as the timepoint drifted\footnote{Mixing monotonic and realtime counters
is a real life issue: GLIBC implementation prior to 2019 strictly applied recommendations found in the STL definition
and used monotonic timers to create \texttt{time checkpoints} and real-time timers to evaluate the exit condition
of \texttt{wait\_for} and \texttt{wait\_until} methods; altering system time --- from the system console for example ---
to be in the past or in the future made applications hang forever}.

The resolution of the system clock is configurable according to the limits imposed
by the clock device driver; this detail is platform dependent and for this reason the driver
offers an I/O control to retrieve the minimum and the maximum values accepted by the hardware counter.
The clock period is dynamic and changes accordingly to the delay of the running threads,
in other words it is tuned to the next timepoint and the value of this timepoint
depends on the workload of the system; if the only running thread is the \texttt{idle thread}
the period will grow constantly until it hits a configured ceiling.
{\myOS} is a \textbf{tickless} kernel.

The \textbf{threads synchronization APIs}
\marginline{\footnotesize{synchronizing threads}}
is a collection of system utilities abstracting
internal kernel structures used to let application alternate each other while accessing
shared resources;
since the \textbf{scheduler} and the threads management are hidden layers, the applications 
have no direct access to each other nor can communicate directly with each
other\footnote{There is no memory protection in {\myOS} so threads can mutually
access private data and variables, while this is strongly discouraged, it is also a
bad practice}.
Applications need to synchronize when they share data, objects, files, or when
they need to observe a timeline --- so they need to synchronize with the
system clock.
{\myOS} implement \textit{mutexes}, \textit{barriers}, \textit{spin locks}, \textit{timers}
and \textit{alarms}.

Additionally, an API for event management is provided to let threads talk to
each other without crossing function calls or sharing global data.
This set of functions \textbf{must be protected} against interruptions;
interrupt routines \textbf{must not interrupt} or alter the running thread state or
stack while the thread is suspending or waking up, as this might change the
readiness state of a thread.
This could result in an erroneous condition for one or more threads leaving the scheduler
in an inconsistent state, or even corrupting a thread's stack.
This sets an important paradigm in the design of {\myOS}: context switch is a \textbf{voluntary}
action taken inside a running thread and executed calling the \texttt{kernel API}.

The \textbf{kernel interface}
\marginline{\footnotesize{calling the kernel functions}}
is the public interface used by application code
to interact with the management of threads, with the system startup and
shutdown, and with system devices.
This interface allow applications to create and destroy a thread, suspend it
and make it sleep for a specific amount of time, and access hardware through
devices.
The true nature of this interface is not hiding thread management, but rather
hiding \textit{scheduling} operations, preventing applications from accessing
any of the \textbf{scheduler} functions.
As a matter of fact, in {\myOS} the scheduler is a hidden layer and is
accessible  only through the \textbf{kernel interface}.
Also hardware and drivers are hidden by this layer, they are accessible by
using the \textit{devices} interface which is an abstraction of the drivers
interface.

\subsection{Layer 2: The libraries}
\textbf{Libraries} are the core of any software infrastructure.
Algorithms and data structures are stored in libraries to share knowledge,
make their implementations common and available to applications, the kernel,
the drivers.
Back in 1976 Niklaus Wirth published a book titled \emph{"Algorithms + Data
Structures = Programs"} creating a milestone for software engineering: data
structures and algorithms are indipendent yet interlocked, as data structures
are owned by the running application, the algorithms are implemented in code;
it is not enough to thoroughly design only one of them.

There are two places for
\marginline{\footnotesize{where libraries are}}
\textbf{libraries} in {\myOS}, the user
applications libraries and the kernel libraries. The difference between the
two is the target they want to accomplish: applications need general
purpose data to solve general purpose problems, so they are offered lists, queues,
binary trees and some strategies to use them at best, like sorting or
searching.
On the other hand the kernel code need specific data to handle drivers, manage
hardware, communications busses, interrupt controllers, but again with common
strategies like enumeration or sorting.

One noticeable detail
\marginline{\footnotesize{the native libc}}
in {\myOS} is the presence of a native \texttt{libc}
library available as original source code in place of one of the several
well-known implementations available in the world.\footnote{At the time of this writing well-known \texttt{libc}
implementations are \texttt{glibc}, \texttt{uClibc}, \texttt{musl} and the
\texttt{C runtime Library (CRT)}} This singularity needs some explanation.
This Operating System is a recreational toy not for production;
porting a library to a new Operating System is a big effort,
expecially considering the \texttt{libc} contents and its portability requirements.
Writing a new dedicated library proved to be effective for performances and size and
required a coding effort equivalent to spending some nights awake with a big mug of coffee.
\footnote{I would not call this burning my midnight oil.}

Kernel libraries offer support for handling the PCI bus, some timing facility,
and interfaces to the file system.
Application libraries offer support for red\&black trees, lists, queues, hashed
lists, memory pools, chunk managers, dynamic vectors, circual buffers, stacks,
and many others.

\subsection{Layer 3: The applications}
\textbf{Applications} are the top level layer in {\myOS} and they are
barely noticeable in the code base, users are expected to add whatever they 
want to this layer, customizing the system experience.
An \textit{application} is not self-contained as it is supposed to be, as the
overall system is a single binary; applications are prevented from accessing
the internals of the kernel and are allowed to use only public APIs as found
in libraries.

Applications code run inside threads, at least one must be spawned at the end
of the kernel initialization by declaring the function \texttt{platform\_run()}
which is responsible for initializing data and threads.
While one binary can run several applications, only one \texttt{platform\_run()}
function can be declared: this means that applications
can be spread over a number of files, but the entry point is common and is
inside \texttt{platform\_run()}.

A number of system functionalities are implemented as applications, leveraging on the fact
that running in the user layer sometimes offers more possibilities than running in the kernel layer
--- libraries, devices, structured I/O can be used to implement a network stack, or a file system,
or a HID interface.

\section{Portability}
The source code of {\myOS} is written in plain C, with a little amount of
assembler to perform hardware init.
The C language is a compulsory choice when portability is a priority, and
along with it the speed of the resulting binary.
Assembler is not required apart for initializations and to access specific
instructions not known to the compiler; it can be handful when it is desidered to 
add hard optimizations to some function or even a whole API, provided that readability
of the code is kept to a level where any developer could read the code,
read the comments, understand the algorithm, the pitfalls,
and the possible C counterpart.\footnote{Assembler by definition is never portable, but can be ported
to different platforms if it is understandable}

Some drivers and some libraries' API have a double implementation, one in plain C and 
one optimized written in assembler; the former is totally portable,
the latter is welcome when performance gets on stage.
The system is not designed to run on 64-bit architectures, writing 32- and 64-bit compliant 
code requires a considerable amount of time expecially if we are considering an
Operating System; software cannot accomodate every difference with a common implementation
when it comes to hardware handling or system management.\footnote{
Writing code able to compile and run on both 32- and 64-bit processors would break
the rule of simplicity; still, you will find 64-bit data structures in the code}

Inlines are extensively adopted in place of macros to gain type checks and 
to prevent funny warnings generated by macro expansion --- it is very easy to loose control
over macros as they are expanded by the preprocessor and then compiled, letting variables
to be declared and added, possibly duplicated, creating redundancies.

The system supports compiling and linking source files written in C++ using
classes and objects, this makes it possibile to write applications with
\textit{Object Oriented Programming}. Unfortunately it was not foreseen to add objects to
the kernel, although possible, it would require a new architecture.

The system can be easily ported to different platforms or can support several
new devices making use of its driver infrastructure and the minimal
\textit{HAL} required to support a board.\\
The \textit{HAL} layer is described in subsection~\ref{subsec:hal} at
page~\pageref{subsec:hal}.

\section{The Build System}
The code base of {\myOS} builds using makefiles and the make tool, no ancillary
tool is required apart for a C-language toolchain.
Any toolchain can be deployed and used to compile and build {\myOS},
at the time of this writing only GCC has been tested.

Output object files are stored in dedicated directories named after the platform 
and CPU configured for the specific build.
Due to this choice it is possible to build several binary images in the same workspace without generating
conflicts.
It is also easier to add platform specific optimizations inside general purpose modules as all makefiles
are aware of the platform and CPU options.
The makefiles \textit{will not} automatically perform a selection of source
code using platform and CPU as a key, this excercise is left to the
developers.

Configuration of the build system is implemented with makefiles, so
there is no need to use other tools rather than \textit{make}
to configure, compile and build {\myOS}.
This choice implies passing the configuration to the toolchain programs as
command line strings, the size of the configuration is then limited by
the length of the command line.
Again, while there are smarter solutions, this is a simple one and will not let the
system become too complex\footnote{A very good alternative to standard Unix makefiles
is the CMake build system, it might be the next build system for {\myOS}}.

\chapter{The System}
\dictum[Dave Cutler]{
Little did I know that I would be fortunate enough to develop several
operating systems in my lifetime; developing one is a rare opportunity for
anyone.
}
\newpage
The kernel, the libraries, the applications, they are all parts of a whole
system. A system with a well defined target: start applications and make them
run in a steady environment.
To this purpose some additional sections are required in {\myOS}, the
boot code to let the main binary run in a preconfigured hardware environment,
the file system to store and retrieve data, the networking subsystem to
connect to the outer world.

\section{The system loader}
All operating systems need to perform a common, early stage operation: load a binary from storage
into memory and then execute it.

This early stage application is commonly known as \textit{loader} or \textit{boot loader}.

A \textit{boot loader} is a very minimal application with limited access to hardware and limited performances,
its solely purpose is to let the hardware run the operating system, which is in charge of performing a full initialization
of the hardware.

Seldom the \textit{boot loader} will not load the operating system, but rather a \textit{second-stage loader} which will load
and execute the operating system; this latter kind of executable is also known as \textit{monitor} and is capable
of breaking and resuming the boot process, perform diagnostics and support early stage debugging.

\medskip
\begin{tikzpicture}[scale=1.25,transform shape,font=\fontsize{9.6pt}{11.52pt}\selectfont]
	\draw[gray,dotted,line width=1.2pt] (4.55,-0.95) rectangle (9.35,0.95);
	\node[circle,draw=Gray,fill=Gray!30!White,text=Black,line width=1.2pt,
		minimum size=2.2cm,align=center] (POWERON) at (0,0) {POWER ON};
	\node[rectangle,draw=NavyBlue,fill=NavyBlue,text=white,line width=1.2pt,
		text width=1.6cm,minimum width=1.6cm,minimum height=1cm,
		inner sep=2pt,align=center] (FBL) at (2.8,0) {First Boot Loader};
	\node[rectangle,draw=Teal,fill=Teal,text=white,line width=1.2pt,
		text width=1.6cm,minimum width=1.6cm,minimum height=1cm,
		inner sep=2pt,align=center] (SBL) at (5.8,0) {Second Boot Loader};
	\node[rectangle,draw=Orange,fill=Chocolate,text=white,line width=1.2pt,
		text width=1.6cm,minimum width=1.6cm,minimum height=1cm,
		inner sep=2pt,align=center] (RUN) at (8.1,0) {Start Operating System};
	\draw[->,thick, dotted] (POWERON.east) -- (FBL.west);
	\draw[->,thick] (FBL.east) -- (SBL.west);
	\draw[->,thick] (SBL.east) -- (RUN.west);
\end{tikzpicture}
\medskip

In {\myOS} the system loader is usually merged on top of the OS header in order to compose
a single binary file, so that the \texttt{SoC} will execute the binary --- either from memory or from ROM --- 
and the execution will start the loader, which will init
the bare minimum hardware, read the content of the remaining binary into memory if the loader did not
provide to do so, and then jump to the entry point of the kernel.
This implies that the hardware platform should be able to execute from a storage device\footnote{An internal or external ROM,
a flash memory device, a hard disk} and then jump to RAM.
This also implies that the processor should be able to bootstrap from storage, or in other words, it should be able to access
storage in hardware with no software intervention.
Some platforms follow a more structured approach where the \textit{boot loader} is an indipendent executable running alone as
a bare metal application.

\section{The boot code}

The boot code belongs to the booting stage and its ultimate goal is to
implement the most straighforward solution to setup the processor, the CPU
peripherals, then load and run the {\myOS} binary.
Boot code should be adaptable\footnote{Adaptable does not imply portable} to any platform and any medium,
and it should perform using the least amount of interactions with the hardware.
It is understood that boot code is essentially processor and platform dependent,
even reducing the assembler parts to the bare minimum there would be a huge
portion of operations which vary with the machine.

There is an elegant and effective solution to the problem of mixing
dependencies with abstract algorithms, the \textbf{Hardware
Abstraction Layer}; in {\myOS} most of the booting process is
implemented into \textbf{HAL} functions, so that a very thin layer of
high-level functions can perform common steps indipendently from the
underlying hardware.

To satisfy these requirements {\myOS} is built as a raw binary, relocated
to start at a specific memory address --- the same that will be used to 
load it and run it\footnote{At the time of this writing there is no support
for COFF, ELF or PE formats}.
An ELF file is built along with the raw binary to satisfy the needs of having
a symbol reference for debugging with \textit{GDB} and to be able to use
all \textit{binutils} tools.
In case the target EABI does not support the ELF format, any other object file
format can be used; just be aware you always need to emit a raw binary file to
load and execute {\myOS}.

Coming to procedures, the boot code implements a detailed setup procedure to have consistent 
booting steps across any platform.
This procedure is described in the following bullets.
Please note the recommendations on how to write and document the code: any detail counts in this stage.

\begin{enumerate}
\item Setup the host processor in order to be able to execute the boot code.
Reserve a register or a memory location to store debugging data like return 
values or failures, it can be very handy even if a debugger is available.
This stage comes so early in the booting process that understanding the
easiest issue might be really hard.
Add comments to explain how this data can be used.
Do not add debugging infrastructures, traces, endless strings for printouts.
\item Access the medium where you stored the {\myOS} binary, and read it into
memory. Document where the binary is put in memory, why it is put there, and
the memory layout during boot. The medium and the code required to
access it are up to you. The rationale here is, split your boot code if it grows beyond
logic; chain booting is not an option in this case.
\item Before jumping to the OS entry point go through this checklist:
\hfill
\begin{tcolorbox}[enhanced, colframe=CadetBlue!80!White, leftrule=3mm, top=5mm, bottom=5mm, halign=center, drop fuzzy shadow]
    \begin{itemize}
    \item All registers must be setup to start {\myOS}: scratch registers should be
    set to zero, memory registers should be set to access the whole addressing
    space; you could also restrict the range of available addresses, but you 
    must ensure the OS is still able to access all the memory and all the hardware.
    \item A small stack will be used by the system init stage, be sure to set
    it up before jumping or you'll crash the machine.
    \item All interrupts \textit{must be disabled} prior to jumping to the
    {\myOS} entry point. 
    \end{itemize}
\end{tcolorbox}
\item Jump to the system entry point, which is labelled as \texttt{DiegOS}.
Be careful not to perform a function call, but simply jump to the label;
calling implies using the stack, and setting a return address somewhere, which
is not expected to happen\footnote{No Operating System jumps to the boot
code after shutting down or after a reset}.
\end{enumerate}

\section{System Initialization}
Initialization is split into a set of function calls, acting on both hardware and
software.

\hfill
\begin{tcolorbox}[enhanced, colframe=YellowGreen!80!White, leftrule=3mm, top=5mm, bottom=5mm, halign=center, drop fuzzy shadow]
\begin{itemize}
\item Initialize the uninitialized data memory space, the .BSS section in
ELF format terminology.
\item Initialize the processor, including power save features, interrupts,
cache memory.
\item Initialize the platform hardware and software; usually you would init
heap and \texttt{SoC} devices here, as well as platform dependent parts 
 like malloc or I/O memory maps.
\textbf{Do not initialize} any driver at this stage.
\item Initialize the debug \texttt{TTY}; depending on your platform, this can be a frame buffer,
a video card, a serial line; keep in mind that this is the system's \texttt{TTY}, and will be used
for debugging purposes only.
\item Initialize the drivers; details about the drivers infrastructure are
explained in section \ref{sub:drivers} at page \pageref{sub:drivers}
\item Initialize the kernel
\item Enable all interrupts
\item Initialize the users' additions, threads added to perform specific tasks
at user's will.
\item Run the kernel, this actually roll the ball and kicks the scheduler.
\end{itemize}
\end{tcolorbox}

\section{The Kernel}
The core functionalities of {\myOS} are enclosed in this part of the software.
The kernel is a single instance, single processor component; it does not
support any kind of memory protection, and it is expecting to be able to access the
whole address space with no restrictions. As mentioned before, there could be restrictions
on the address space, but the OS must still be able to drive all the hardware and must
be able to access all the memory address space.
The kernel does not support threads contention or access privilege, so any portion of
the code (and consequently any thread) is able to peek or poke any memory location or 
access any hardware device.
This is a questionable choice, although it is \textit{"nothing the God of Information wouldn't let you in 
heaven for"}\footnote{I let the reader guess the quote and the movie}.
All in all, a general purpose Operating System would not even mention memory protection as an option; 
again, if you need the real thing, go get it. A small bare-metal system would require lots
of resources to assure a decent level of security when implementing memory protection,
and this is not a goal for {\myOS}.

There are a number of services offered by the kernel to support applications with basic operations,
in a away that it is possible to keep a separation between processes in term of resources and code.
These services offer synchronization with time, between processes, resource contention facilities and
input/output management.
A schema of the structure of these services is sketched in figure \ref{fig:tikz-thekernel} at page \pageref{fig:tikz-thekernel}. 

\begin{figure}
	\centering
	\resizebox{\textwidth}{!}{%
	\begin{tikzpicture}[
	    x=1cm,
	    y=1cm,
	    >=Latex,
	    line width=0.7pt,
	    font=\small,
	    block/.style={draw, fill=gray!20, align=center, minimum width=2.2cm, minimum height=1.2cm},
	    redblock/.style={draw=red!70!black, fill=red!15, align=center, minimum width=2.4cm, minimum height=1.4cm},
	    bluecircle/.style={draw=blue!70!black, fill=blue!20, circle, minimum size=2.1cm, align=center},
	    schedulercircle/.style={draw=orange!80!black, fill=orange!20, circle, minimum size=2.1cm, align=center},
	    aquamarinecircle/.style={draw=Aquamarine!70!black, fill=Aquamarine!45, circle, minimum size=2.1cm, align=center},
	    kernelcircle/.style={draw=black, fill=white, circle, minimum size=5.6cm},
	    groupbox/.style={draw=black, fill=none, inner sep=8pt},
	    filebox/.style={draw, dashed, minimum width=4.4cm, minimum height=3.0cm, align=center},
	    netbox/.style={draw, minimum width=4.8cm, minimum height=3.2cm, align=center},
	    tiny/.style={draw, fill=gray!20, minimum width=1.8cm, minimum height=1cm, align=center}
	]

	% Kernel shell
	\node[kernelcircle] (kernel) at (0,0) {};
	\node[bluecircle] (core) at (-0.9,-0.6) {Kernel};
	\node[schedulercircle] (sched) at (0.9,-0.6) {Scheduler};
	\node[aquamarinecircle] (clock) at (-0.1,0.95) {System\\Clock\\Manager};

	% Top-level components connected to the kernel
	\node[redblock] (boot) at (-4.4,4.0) {Boot\\Code};
	\node[redblock] (pdc) at (4.2,4.0) {Hardware\\Abstraction\\Layer};

	% Lower-left service blocks
	\node[block] (alarms) at (-5.0,-4.7) {Alarms\\Manager};
	\node[block] (events) at (-2.5,-4.7) {Events\\Manager};
	\node[block] (io) at (-3.65,-6.2) {I/O\\Manager};
	\node[groupbox, fit=(alarms)(events)(io)] (servicebox) {};
	\node[below=2pt of servicebox.south] {Events};

	% Right-hand concurrency blocks
	\node[block] (mutex) at (2.65,-4.7) {Mutexes};
	\node[block] (spin) at (2.65,-6.2) {Spin Locks};
	\node[block] (timers) at (5.15,-4.7) {Timers};
	\node[block] (barriers) at (5.15,-6.2) {Barriers};
	\node[groupbox, fit=(mutex)(spin)(timers)(barriers)] (concurrencybox) {};
	\node[below=2pt of concurrencybox.south] {Concurrency};

	% Arrows from external components
	\draw[-{Latex[length=2.5mm]}] (boot) -- (kernel);
	\draw[-{Latex[length=2.5mm]}] (pdc) -- (kernel);
	\draw[-{Latex[length=2.5mm]}] (servicebox) -- (kernel.south west);
	\draw[-{Latex[length=2.5mm]}] (concurrencybox) -- (kernel.south east);

	\end{tikzpicture}
	}
	\caption{The kernel}
	\label{fig:tikz-thekernel}
\end{figure}

\section{Kernel Services}
The kernel offers the following services: a thread scheduler, a thread management API, 
an alarm API, events queueing and dispatching, mutexes, barriers, I/O waiting and spin locks.\\
Alarms, Events and Mutexes are all managed services, meaning that an owning thread
can create, destroy, or reconfigure them at any time.

\begin{description}
\item[The Scheduler] the {\myOS} scheduler is as simple as possible, and is implemented as a FIFO
queue with several priority levels. Pre-emption is not supported, and threads must co-operate in
 order to give up CPU time when computing is not required.

Threads can be suspended, halted, delayed, and removed\footnote{Removal means the
thread will be killed.} but at least one thread must be ready to run at any time
\footnote{This is the Idle thread, which cannot be killed.}.

\definecolor{runningpurple}{HTML}{DA77F2}
\definecolor{runningpurpleedge}{HTML}{9C36B5}
\definecolor{idlefill}{HTML}{B2F2BB}
\definecolor{idleedge}{HTML}{2F9E44}
\definecolor{normalfill}{HTML}{A5D8FF}
\definecolor{normaledge}{HTML}{1971C2}
\definecolor{highfill}{HTML}{FFF9DB}
\definecolor{highedge}{HTML}{FFD700}
\definecolor{realtimefill}{HTML}{FFC9C9}
\definecolor{realtimeedge}{HTML}{E03131}
\definecolor{ink}{HTML}{1E1E1E}

\begin{figure}
    \centering
\resizebox{\textwidth}{!}{%
\begin{tikzpicture}[
    x=0.02cm,
    y=-0.02cm,
    line width=0.04cm,
    font=\small,
    >=Latex
]
\begin{scope}[shift={(-641,-178)}]
\path[use as bounding box] (641,178) rectangle (1258,1122);

% Running thread and its priority-state transitions
\draw[draw=runningpurpleedge, fill=runningpurple]
    (961,251) ellipse [x radius=1.5cm, y radius=1.46cm];
\node[align=center, text=Black] at (961,251) {Running\\Thread};

\draw[<->, draw=ink] (711.4,376) -- (889.6,285.7);
\draw[<->, draw=ink] (873.2,377) -- (911.7,313.7);
\draw[<->, draw=ink] (1036,382) -- (999,320);
\draw[<->, draw=ink] (1198.9,379) -- (1024.3,291.9);

\draw[->, draw=ink] (976,436) -- (932,436);
\draw[->, draw=ink] (811,434) -- (768,434);
\draw[->, draw=ink] (1140,439) -- (1096,439);

\draw[draw=idleedge, fill=idlefill, rounded corners=3pt] (645,380) rectangle (763,495);
\node[text=idleedge] at (704,437.5) {IDLE};
\draw[draw=normaledge, fill=normalfill, rounded corners=3pt] (808.5,379.5) rectangle (925.5,494.5);
\node[text=normaledge] at (867,437) {NORMAL};
\draw[draw=highedge, fill=highfill, rounded corners=3pt] (975.5,382.5) rectangle (1093.5,497.5);
\node[text=highedge] at (1034.5,440) {HIGH};
\draw[draw=realtimeedge, fill=realtimefill, rounded corners=3pt] (1139.5,383.5) rectangle (1256.5,498.5);
\node[text=realtimeedge] at (1198,441) {REALTIME};

% Idle-priority thread queue
\draw[->, draw=ink, line width=0.02cm] (704,498) -- (704,525);
\draw[->, draw=ink, line width=0.02cm] (702,646) -- (702,675);
\draw[->, draw=ink, line width=0.02cm] (702,798) -- (701,826);
\draw[draw=idleedge, rounded corners=3pt] (643,527.5) rectangle (761,642.5);
\node[text=idleedge] at (702,585) {Thread 1};
\draw[draw=idleedge, rounded corners=3pt] (642,679.5) rectangle (760,794.5);
\node[text=idleedge] at (701,737) {Thread 2};
\draw[draw=idleedge, rounded corners=3pt] (642,829.5) rectangle (760,944.5);
\node[text=idleedge] at (701,887) {Thread 3};

% Normal-priority thread queue
\draw[->, draw=ink, line width=0.02cm] (865,499) -- (866,525);
\draw[->, draw=ink, line width=0.02cm] (865,645) -- (865,676);
\draw[->, draw=ink, line width=0.02cm] (866,795) -- (865,826);
\draw[draw=normaledge, rounded corners=3pt] (807.5,525.5) rectangle (924.5,640.5);
\node[text=normaledge] at (866,583) {Thread 1};
\draw[draw=normaledge, rounded corners=3pt] (805.5,678.5) rectangle (922.5,793.5);
\node[text=normaledge] at (864,736) {Thread 2};
\draw[draw=normaledge, rounded corners=3pt] (802.5,828.5) rectangle (919.5,943.5);
\node[text=normaledge] at (861,886) {Thread 3};

% High-priority thread queue
\draw[->, draw=ink, line width=0.02cm] (1037,501) -- (1038,527);
\draw[->, draw=ink, line width=0.02cm] (1040,644) -- (1040,675);
\draw[draw=ink, line width=0.08cm] (1006.5,822.5) -- (1067.5,821.5);
\draw[draw=highedge, rounded corners=3pt] (978,525.5) rectangle (1096,640.5);
\node[text=highedge] at (1036.5,583) {Thread 1};
\draw[draw=highedge, rounded corners=3pt] (977,677.5) rectangle (1095,792.5);
\node[text=highedge] at (1036,735) {Thread 2};
\draw[draw=highedge, rounded corners=3pt] (980,849.5) rectangle (1098,964.5);
\node[text=highedge] at (1039,907) {Thread P};

% Real-time thread queue
\draw[->, draw=ink, line width=0.02cm] (1202,503) -- (1203,529);
\draw[draw=ink, line width=0.08cm] (1169.5,676.5) -- (1230.5,675.5);
\draw[draw=realtimeedge, rounded corners=3pt] (1139.5,526.5) rectangle (1256.5,641.5);
\node[text=realtimeedge] at (1198,584) {Thread 1};
\draw[draw=realtimeedge, rounded corners=3pt] (1141.5,707.5) rectangle (1258.5,822.5);
\node[text=realtimeedge] at (1200,765) {Thread Q};

% Additional thread slots
\draw[draw=ink, line width=0.08cm] (670,971) -- (731,970);
\draw[draw=ink, line width=0.08cm] (834.5,969.5) -- (895.5,968.5);
\draw[draw=idleedge, rounded corners=3pt] (641,1004.5) rectangle (759,1119.5);
\node[text=idleedge] at (700,1062) {Thread M};
\draw[draw=normaledge, rounded corners=3pt] (803.5,1006.5) rectangle (920.5,1121.5);
\node[text=normaledge] at (862,1064) {Thread N};
\end{scope}
\end{tikzpicture}
}

\caption{The cooperative scheduler}
\label{fig:tikz-thescheduler}
\end{figure}

The scheduler is designed to manage the threads' state machine and select 
the running thread by using a priority queue schema. There are four
priorities in {\myOS}, namely \textit{Idle Priority}, \textit{Normal Priority}, 
\textit{High Priority} and \textit{Realtime Priority}.

When a thread is in the ready queue, it will not substitute the running
one, not even if the ready thread's priority is higher; pre-emption is \textbf{not
supported} at all.
The running thread must give up by suspending or waiting on some shared resource to let
the scheduler select another running thread.

The scheduler will loop the highest priority queue with at least a ready thread until the
list is void; the list will be empty when all the threads with that priority have suspended or have
been killed.
This implies that all threads in a specific ready queue will run at least once before being evicted
from the list itself\footnote{Eviction from a ready queue is voluntary and can be performed by waiting or sleeping or dieing}.
Lists are browsed in decreasing priority order from \texttt{Realtime Priority} to \texttt{Idle Priority} where at least one ever-running
thread is expected to exists\footnote{The ever-running thread is the idle thread}.

Whatever the thread's priority might be, a thread owns the whole system until suspending or waiting: care must be taken
in writing application code processing data to avoid hogging the CPU or - definitively worser - starving other
threads.
This is a design principle in {\myOS}: threads have a \textit{Priority} to be selected for execution
but they all have the same execution \textit{privilege}\footnote{Same execution privilege includes
equal access to privileged instructions, I/O ports, memory, registers, counters}.

Bear in mind that a higher priority thread will never give up for a lower priority
thread; to avoid CPU hogging you must be sure that all threads will either sleep or 
wait or delay, and they should perform their processing as fast as possible;
if you cannot complete the required actions in a decent amount of time, you must split
operations and run them in chunks, suspending the thread between runs;
failing to do so will starve all threads exept the running one.
Even if this might be a desired behavior, it is really dangerous and might generate unwanted effects
on accounting or performance monitoring.
Generally speaking it is a best practice to split a task load into slices and run slices individually,
yelding to other threads at every slice.

\item[Threads] a thread is a new independent code path starting from a
function (used as entry point) that can be alternated with other code paths.
Whenever a thread hands over to another thread, the kernel has to store the
CPU's context into the running thread, and restore the context loading it from
the handed over thread.
Context is defined as the set of registers and flags that differ from path to
path in the code.\\
Along with the processor context a distinctive feature of a thread is its \textbf{stack}.
The stack pointer (which is usually part of the processor context) must be saved and restored
and stack memory areas \textbf{must not} be corrupted at any time, or the distinctive backtrace of 
code paths would break leading to memory corruptions or crashes.

\begin{figure}
\centering
\begin{tikzpicture}[->,  outer sep=auto, inner sep=1em, thick, node distance=4cm]
\node [circle,fill=gray!30,draw] (IDLE) {IDLE};
\node [circle,fill=LimeGreen,draw] (READY) [below of=IDLE] {READY};
\node [circle,fill=Orange,draw] (RUN) [below of=READY] {RUN};
\node [circle,fill=gray!60,draw] (KILL) [below  of=RUN] {KILL};
\node [circle,fill=RoyalBlue,draw] (SLEEP) [left of=KILL] {SLEEP};
\node [circle,fill=RoyalBlue,draw] (WAIT) [right of=KILL] {WAIT};
\path
(IDLE)     edge [dashed, loop above] node [right] {On creation} (IDLE)
(IDLE)     edge [dashed] node [right] {After creation} (READY)
(READY)  edge node [right] {Scheduled} (RUN)
(RUN)      edge [bend left=60]  node [left]{suspend} (READY)
(RUN)      edge node [left]{sleep} (SLEEP)
                edge node [left] {wait} (WAIT)
               edge [dashed] node [left] {killed} (KILL)
(SLEEP)   edge [bend left=60] node [left] {wake up} (READY)
(WAIT)   edge [bend right=60] node [right] {wake up} (READY);
\end{tikzpicture}

\caption{Thread State Machine }
\label{fig:tikz-threadstatemachine}
\end{figure}

Threads can be in four stable states, namely \texttt{RUN, WAIT, SLEEP,  READY}, and two transient states
 \texttt{IDLE} and \texttt{KILL}.
The latter are transient by definition as they designate a thread which has been created but never run (\texttt{IDLE})
and a thread scheduled to be killed (\texttt{KILL})
\footnote{A thread is not killed upon call to a specific kernel interface, it is inserted in a temporary queue for later disposal}.\\
Threads can be promoted from the \texttt{READY} state to the \texttt{RUN} state  by the scheduler internals only;  the reverse action is performed by \textit{suspension} function calls.\\
Kernel APIs can change the \texttt{RUN} thread state to \texttt{WAIT} state by means of \textit{sleeping} functions or \textit{suspension points}
\footnote{Kernel suspension points are mutexes, barriers, I/O waits}.

\item[Alarms] are a fundamental feature to let threads sleep and wake up at a
defined time. 
Alarms are events queued to a thread with a specific class ID; handling alarms is just
like handling events. Alarms can be one-shot or recursive, and can be started or
halted anytime.
\item[Events] are extensible system messages, used to signal some interesting news to a thread. 
Events are organized as queues (FIFO), and processed one by one in
chronological order.
Memory management for events is optional,  both the producer and the
consumer must take care of allocating memory when creating an event message,
and freeing it once the event has been processed. The producer can create events
with no memory-free handler, in this case the consumer will not release the event storage;
if a memory-free handler is set, the consumer will invoke it passing a pointer to the event
as a parameter to the handler.

\begin{warning}
\textbf{Important:} the memory release handler (memory-free handler) is executed in
the \textit{consumer} thread's context, possibly out of control of the producer's thread;
if events are created on the heap a synchronization mechanism must be put in place to
avoid corrupting data structures.
\end{warning}

\item[Mutexes] are the simplest form of semaphore. A mutex is initially
unlocked; the first thread trying to acquire a lock on the mutex will actually get the lock
and proceed with execution, entering a critical code section and running as long as it needs to.
Unless the mutex is unlocked, no other thread is allowed to enter the same
critical code section; other threads trying to acquire a lock on the same mutex will sit in a
waiting queue. 
Once the first thread has released (unlocked) the mutex, the threads waiting for 
a lock will be run one at a time, until no waiting threads are left; every thread entering the
critical section is guaranteed to have acquired a lock on the mutex, which means, 
code after the lock acquisition function is guaranteed to run helding a lock on the mutex.
Once no thread is waiting for a lock on the mutex, the mutex is left unlocked.
\item[Spin locks] are a running  atomic delay to acquire a lock on a integer used as a
barrie. Spin locks can be locked or unlocked.
A thread trying to acquire a spin lock will enter an infinite loop until the lock is 
gained, so if the spin lock is originally locked, the thread will suspend and enter
the ready queue continuosly.
Good use of spin locks is to protect very short critical sections or hardware accesses,
where the longer code path required to manage mutexes might be not as efficient.
\item[I/O waits] are wait queues, useful to suspend the calling thread execution till 
data is available, this usually happens while waiting for input from a hardware device.
The suspended thread is removed from the ready queue and parked into a wait queue, multiple
threads can sleep on the same wait queue, and can be awakened in groups.
Threads are moved from wait queues to ready queues in a FIFO fashion, and it is not possible
to priorize or sort this sequence.
\item[Barriers] are a form of cumulative mutual exclusion, useful to suspend execution of several
threads waiting for a single shared resource to be available.
Barriers are \textit{closed} upon creation, and can be opened or closed by any thread; an open barrier will
resume \textbf{all} the waiting threads at once moving them into their respective \textit{ready queues}.
Optionally a barrier can automatically close after opening, creating a step-lock software gate.
\end{description}

Whenever a thread is put in \texttt{WAIT} state it is inserted in a shared \textit{wait queue} where
priority does not matter.
Threads in this queue does not run and cannot be elected to execute; the scheduler will periodically
scan the queue and wake up all eligible threads to be moved into their respective \textit{ready queues}.
Threads waiting on I/O or waiting for an infinite amount of time are woken up by other primitives of the
kernel, their state is changed to \textit{READY} but they are not removed from the \textit{wait queue},
this task is left to the scheduler.

\section{Drivers}
\label{sub:drivers}
Drivers are a software interface conceived to establish a communication with hardware.
They enable software to start and stop hardware operations, configure it, and perform input and output.
Common operations performed by drivers include enabling and handling interrupts: the pertinent code must
be carefully designed to avoid conflicts with the system-wide interrupt handling.
%{\myOS} protects critical sections by enabling and disabling interrupts, and these operations \textbf{are}
%designed to be reentrant.
%The system-wide interrupt handler will invoke the registered callback with the corresponding interrupt \textbf{off},
%and will enable the interrupt back again once the servicing routine returns.

\subsection{A common interface}
Hardware comes under different forms, but certain operations apply to all of them.
A device need to be initialized --- configured --- in order to work as expected and to perform
specific routine with specific options.
As an example, a disk drive controller can use DMA transfers when reading or writing,
and it will need a configured memory address to tranfer data from the disk drive to memory and reverse.
The same controller might accept to work disabling DMA transfers, in this case a different configuration
is required to handle data transfers and the driver will need a different timing strategy to make the hardware
perform decently
\footnote{DMA memory transfers are completely managed by hardware, programmed transfers require software to move data
from the device to the central memory}.

In {\myOS} there is a common interface valid for all drivers, and a set of rules used to define it.
All functions returning a status will return an \texttt{int} value which specifies if the operation
was successful or it encountered a problem, in the latter case the type of error is reported
using the values found in Unix or Linux and set into the \texttt{errno} variable;
this error reporting schema is far less than optimal but simple enough for a simple Operating System.
All functions returning a size will return an \texttt{unsigned int} value.
Drivers are multi-instance, as a consequence all function prototypes must have an \texttt{unsigned int} parameter
to identify the instance number.
\vfill
\begin{warning}
The design of the drivers' interface is going to change in the future replacing parametrics tuple (driver, unit)
with a more consistent context-based approach.
\end{warning}
\vfill
Let's now break down the drivers' common header and its structures, you can have a look at the
file in listing \ref{list:drivers_common}.

\begin{description}
\item[name] is character string storing the driver's name, this is a quite short string and should
resume the name or purpose of the driver, usually the name of the device to be driven.
\item[init()] function can be used to initialize the driver for a specific instance
\footnote{An instance matches a hardware device, a single driver can handle multiple devices.}
but it \textbf{will not} make the device run: the function contract is to setup the hardware, not make it run.
\item[start(),stop()] can run the hardware activities or stop them, this includes enabling or disabling
interrupts, DMA transfers, I/O operations, sensors polling, and so on.
It is assumed that alternating calls to \texttt{start()} and \texttt{stop()} will make the hardware device either run or not run,
but will not change its configuration or power state.
These two functions --- along with the interrupt routines --- are the most dangerous part of a driver, they deal with interrupts and
buffers, they may leak memory if not properly coded, 
\item[done()] can be used to terminate the driver and shut down the hardware device.
It is coupled with \texttt{init()} just like \texttt{start()} and \texttt{stop()} functions.
Once \texttt{done()} has been invoked it is required to call \texttt{init()}
to reset the device to defaults and reconfigure it.
\item[status()] can be called at any time to report the driver's internal state, the ongoing operation, and the driver's specialization.
\footnote{Drivers have different extensions for different kind of hardware devices: block devices, character devices, graphics, networking}.
\end{description}

The common header and its common functions enable the kernel to initialize, start, stop, terminate and get a report 
of the status of any driver indipendently from its specialization;
all drivers can be managed by the kernel in the same way,
the devices linked to the drivers will inherit their names from the drivers names and
will be able to operate any driver using one semantic only.

This simplifies the structure of the kernel and let the developers add support to their favourite hardware
without the burden of adding initialization sections or functions dedicated to the new driver.
The plug-and-play design of the driver interface is a nice example of the \textit{KISS} principle
and it is one of the most common pattern found in an Operating system
\footnote{This would be obvious in an Operating System coded with an object-oriented language}.

\scriptsize
\inputminted{C}{../include/diegos/drivers_common.h}
\captionof{listing}{drivers\_common.h}
\label{list:drivers_common}
\normalsize
\hfill

\subsection{Character drivers}
A character driver can manage any hardware or virtual device which is handling data in the form
of a stream of bytes --- characters if you'd like.

A character driver expects to input a stream of bytes from the device, and to output a stream of bytes
to the device; the length of the stream is not predefined and can vary from 1 to anything, the stream
might not even have an end and lasts fot the entire running time of the device.
The basic principle of this type of driver is the stream: input and output comes in the form of a sequence
of bytes.
Examples of this type of driver are the serial lines drivers (RS232 or IEEE 485) or the keyboard.

The character driver is the most simple driver and its extended interface is composed by 4 function calls,
namely read and write with a single buffer, and read and write with multiple buffers.

You can have a look at the interface in listing \ref{list:char_drivers}

\scriptsize
\inputminted{C}{../include/diegos/char_drivers.h}
\captionof{listing}{char\_drivers.h}
\label{list:char_drivers}
\normalsize

\subsection{Network drivers}
A network driver can handle devices with networking capabilities, sending or receiving data to remote peers
in chunks of data and with a well-defined set of rules for addressing, encoding and decoding.
A network driver looks similar to a character driver, they both send and receive data streams and --- apart for
the structured data required by a network card --- they are agnostic about the contents of the data.

A character driver sorts data in bytes, a network driver sorts data in packets or frames: the input stream
is organized as a queue of chunks instead of a queue of bytes.

A network driver need additional informations to properly run its device, like the addresses of the sender
and the receiver, the line status (which is \textbf{not} the device's status!), the encoding to be applied to the stream,
the possible multiplexing of senders and receivers.

The extended interface comprises numerical informations like the MTU
\footnote{Maximum Transmission Unit} of the network card, addresses to be used for multicast or unicast transmissions,
timers for a State Machine in case the line protocol requires it, and the usual single - and multiple read and write functions.

You can have a look at the interface in listing \ref{list:net_drivers}

\scriptsize
\inputminted{C}{../include/diegos/net_drivers.h}
\captionof{listing}{net\_drivers.h}
\label{list:net_drivers}
\normalsize

\subsection{Block drivers}
A block driver is quite different when compared with other drivers, not because of its extended interface --- which is very
similar to that of a character driver --- but for the peculiarities of the devices it manages.

Block drivers are structured to support storage in the form of Hard Disks, Solid State Disks, NOR, NAND, eMMC, SD, floppy disks,
and anything else might exist.
All these devices share a common property: they offer pre-sized storage unit called blocks, sectors or pages which are accessed
as a whole.
As an example, a Hard Disk internal memory is usually fractioned in 512-bytes sectors; it is not possible to read or write a part
of a sector, but only the whole 512 bytes. If the driver is submitted with an amount of data which is not multiple of 512 bytes
the data must be padded
\footnote{Additional spare data is appended to the buffer until its size is a multiple of the sector size}
to fill an integral number of sectors; equally if the driver is requested to read an amount of data which is not multiple of 512 bytes
the driver needs to read additional, unwanted data in an internal buffer, then partially copy the last sector to the requester.

Block drivers also need to know \textit{where} data are to be read or written; as a matter of fact, a storage device has a finite capacity
and cannot save an infinite amount of data. A serial line (character driver) can stream the english alphabet for eons until the end of the world
and won't ever need to know when the stream started or when it started.
A block device must be addressed before a write operation by specifing the first sector to be used; the data stream will be saved starting with
the selected sector, proceeding with the next in order until the last sector is reached, at which point the operation will roll to the very first
sector on disk and proceed until the previously selected starting sector is reached; when this second condition is met the disk has been completely
written and no more free storage is available.
The driver is allowed to overwrite previously written sectors, and can proceed in this action for eons until the end of the world, but always
being aware of which sector is being written and which was the first one selected to start operations.

Rolling over the last sector, padding on write and partial copying on read are not present in other type of drivers.

You can have a look at the interface in listing \ref{list:block_drivers}

\scriptsize
\inputminted{C}{../include/diegos/block_drivers.h}
\captionof{listing}{block\_drivers.h}
\label{list:block_drivers}
\normalsize


\section{The File System}

\section{Libraries}

Every operating  environment offers a set of libraries to be used
at application level.
Applications make use of widely known data structures and algorithms
to build even more complex algorithms.
When an Operating System is built as a single process, running application
level threads, the availability of such APIs enable developers to code
highly orthogonal systems, leveraging on the same libraries for drivers,
front ends, interfaces, databases, applications.

{\myOS} offers a pletora of data structures
\section{Networking}

\section{Drivers interface}

\chapter{The File System architecture}
\dictum[John McKittrick]{
There's no way that a high school punk can put a dime into a telephone and break into our system!
He's got to be working with somebody else. He's got to be!
}
\newpage
\lipsum

\chapter{The Networking infrastructure}
\dictum[Maverick and Goose]{
I feel the need... the need for speed!
}
\newpage
\lipsum

\chapter{The System Console}
\dictum[Andrew S. Tanenbaum]{
I can type faster than I can point. And my mother told me that pointing is impolite.
}
\newpage
\lipsum

\chapter{The Debug Infrastructure}
\dictum[Travis Penn]{
Assumption is the mother of all fuck-ups!
}
\newpage
\lipsum

\chapter{The Kernel Implementation}
\dictum[Morpheus]{
This is your last chance.\\
After this, there is no turning back.\\
You take the blue pill --- the story ends,
you wake up in your bed and believe whatever
you want to believe.\\
You take the red pill --- you stay in Wonderland
and I show you how deep the rabbit hole goes.
}
\newpage
\subsection{Recommendations and other suggestions}
The core functionalities are very simple: spawn a number of execution threads, run all
of them one by one, interrupts \textbf{cannot and must not} invoke a context switch,
synchronize the execution by suspension, share resources with mutexes or spinlocks.

A few words about mutexes and spinlocks and their pros and cons: mutexes have a more complex
code path and disable or enable interrupts to avoid contention within the scheduler, while
spinlocks are very light and simple.
On reverse, a thread waiting to acquire a lock on a mutex will enter a wait queue, thus it will
sit and wait (and shut up!) until its time has come to run and gain a lock; a thread waiting to
acquire a lock on a spinlock will continuosly suspend and thus will be rescheduled until its
time has come to gain a lock.

The balance is, more complexity but lesser CPU usage, opposed to extreme simplicity but a busy CPU.
My suggestion is to use spinlocks for very fast and short critical sections, whereas managing mutexes
would be worser and slower; make use of mutexes if you need heavy computations inside the critical 
section, as all other contending thread will be waiting and won't be scheduled at all, freeing up the 
CPU and letting the locking thread complete its tasks as fast as possible.

\end{document}
